\documentclass[10pt,a4paper]{amsart}
\usepackage{amsmath,amssymb,amsthm,graphicx}
% Retained preamble for the Stacks Project chapter {CH} (shared project preamble).
% The chapter file pinned-original-source.tex only carries \input{{preamble}} on its first line,
% so no preamble text lives inside the chapter itself.  The Stacks Project uses one single shared
% preamble for every chapter: the snapshot stages it three times (sources/stacks-chow/preamble.tex,
% sources/stacks-etale-cohomology/preamble.tex and focused-seeds/AGP/sources/stacks-snapshot-20261011/
% preamble.tex) and those copies are identical in size.  The block below is copied mechanically,
% line for line, from sources/stacks-chow/preamble.tex lines 14-18 and 154-217: the Xy-pic setup and
% the native theorem declarations together with the mathematical macros.  Nothing was added, removed
% or redefined; the wrapper class and the amsmath/amssymb/amsthm/graphicx base packages are supplied
% by the assembler.
% For commutative diagrams we use Xy-pic
% Chapter: stacks-cohomology.tex -- Cohomology of Sheaves.
\usepackage[all]{xy}

% We use 2cell for 2-commutative diagrams.
\xyoption{2cell}

% Theorem environments.
%
\theoremstyle{plain}
\newtheorem{theorem}[subsection]{Theorem}
\newtheorem{proposition}[subsection]{Proposition}
\newtheorem{lemma}[subsection]{Lemma}

\theoremstyle{definition}
\newtheorem{definition}[subsection]{Definition}
\newtheorem{example}[subsection]{Example}
\newtheorem{exercise}[subsection]{Exercise}
\newtheorem{situation}[subsection]{Situation}

\theoremstyle{remark}
\newtheorem{remark}[subsection]{Remark}
\newtheorem{remarks}[subsection]{Remarks}

\numberwithin{equation}{subsection}

% Macros
%
\def\lim{\mathop{\mathrm{lim}}\nolimits}
\def\colim{\mathop{\mathrm{colim}}\nolimits}
\def\Spec{\mathop{\mathrm{Spec}}}
\def\Hom{\mathop{\mathrm{Hom}}\nolimits}
\def\Ext{\mathop{\mathrm{Ext}}\nolimits}
\def\SheafHom{\mathop{\mathcal{H}\!\mathit{om}}\nolimits}
\def\SheafExt{\mathop{\mathcal{E}\!\mathit{xt}}\nolimits}
\def\Sch{\mathit{Sch}}
\def\Mor{\mathop{\mathrm{Mor}}\nolimits}
\def\Ob{\mathop{\mathrm{Ob}}\nolimits}
\def\Sh{\mathop{\mathit{Sh}}\nolimits}
\def\NL{\mathop{N\!L}\nolimits}
\def\CH{\mathop{\mathrm{CH}}\nolimits}
\def\proetale{{pro\text{-}\acute{e}tale}}
\def\etale{{\acute{e}tale}}
\def\QCoh{\mathit{QCoh}}
\def\Ker{\mathop{\mathrm{Ker}}}
\def\Im{\mathop{\mathrm{Im}}}
\def\Coker{\mathop{\mathrm{Coker}}}
\def\Coim{\mathop{\mathrm{Coim}}}

% Boxtimes
%
\DeclareMathSymbol{\boxtimes}{\mathbin}{AMSa}{"02}

%
% Macros for moduli stacks/spaces
%
\def\QCohstack{\mathcal{QC}\!\mathit{oh}}
\def\Cohstack{\mathcal{C}\!\mathit{oh}}
\def\Spacesstack{\mathcal{S}\!\mathit{paces}}
\def\Quotfunctor{\mathrm{Quot}}
\def\Hilbfunctor{\mathrm{Hilb}}
\def\Curvesstack{\mathcal{C}\!\mathit{urves}}
\def\Polarizedstack{\mathcal{P}\!\mathit{olarized}}
\def\Complexesstack{\mathcal{C}\!\mathit{omplexes}}
% \Pic is the operator that assigns to X its picard group, usage \Pic(X)
% \Picardstack_{X/B} denotes the Picard stack of X over B
% \Picardfunctor_{X/B} denotes the Picard functor of X over B
\def\Pic{\mathop{\mathrm{Pic}}\nolimits}
\def\Picardstack{\mathcal{P}\!\mathit{ic}}
\def\Picardfunctor{\mathrm{Pic}}
\def\Deformationcategory{\mathcal{D}\!\mathit{ef}}

% --- appended 2026-10-12 (Chief): the project own definitions of the reference/slogan/history
% display environments, copied mechanically from the shared Stacks preamble lines 8-12
% (\usepackage{verbatim} plus the three \newenvironment lines, which the chapter preamble keeps
% verbatim from the project preamble). Nothing added, removed or redefined: these environments
% carry no mathematical content and the project itself typesets them as comments. Without these
% three lines a counted statement whose own text uses them fails with Environment slogan undefined.
\usepackage{verbatim}
\newenvironment{reference}{\comment}{\endcomment}
\newenvironment{slogan}{\comment}{\endcomment}
\newenvironment{history}{\comment}{\endcomment}
\usepackage{cleveref}
\crefname{originallemma}{original lemma}{original lemmas}
\Crefname{originallemma}{Original lemma}{Original lemmas}
\makeatletter
\newlabel{lemma-internal-hom-evaluate}{{0A8U}{source}{}{sourcepointer-1}{}}
\newlabel{lemma-internal-hom-evaluate@cref}{{[originallemma][0][]0A8U}{[1][0][]source}{}{}{}}
\makeatother
\crefname{originallemma}{original lemma}{original lemmas}
\Crefname{originallemma}{Original lemma}{Original lemmas}
\makeatletter
\newlabel{lemma-internal-hom-evaluate-isom}{{0G40}{source}{}{sourcepointer-2}{}}
\newlabel{lemma-internal-hom-evaluate-isom@cref}{{[originallemma][0][]0G40}{[1][0][]source}{}{}{}}
\makeatother
\crefname{originallemma}{original lemma}{original lemmas}
\Crefname{originallemma}{Original lemma}{Original lemmas}
\makeatletter
\newlabel{lemma-restriction-RHom-to-U}{{08DL}{source}{}{sourcepointer-3}{}}
\newlabel{lemma-restriction-RHom-to-U@cref}{{[originallemma][0][]08DL}{[1][0][]source}{}{}{}}
\makeatother
\crefname{originallemma}{original lemma}{original lemmas}
\Crefname{originallemma}{Original lemma}{Original lemmas}
\makeatletter
\newlabel{lemma-dual}{{08I1}{source}{}{sourcepointer-4}{}}
\newlabel{lemma-dual@cref}{{[originallemma][0][]08I1}{[1][0][]source}{}{}{}}
\makeatother
\crefname{originalequation}{original equation}{original equations}
\Crefname{originalequation}{Original equation}{Original equations}
\makeatletter
\newlabel{equation-eval}{{08I2}{source}{}{sourcepointer-5}{}}
\newlabel{equation-eval@cref}{{[originalequation][0][]08I2}{[1][0][]source}{}{}{}}
\makeatother
\crefname{originallemma}{original lemma}{original lemmas}
\Crefname{originallemma}{Original lemma}{Original lemmas}
\makeatletter
\newlabel{lemma-Rhom-strictly-perfect}{{08DM}{source}{}{sourcepointer-6}{}}
\newlabel{lemma-Rhom-strictly-perfect@cref}{{[originallemma][0][]08DM}{[1][0][]source}{}{}{}}
\makeatother
\title{The dual of a perfect complex is perfect and double dualization is the identity}
\author{The Stacks Project Authors}
\date{stacks-project:cohomology.tex:local-snapshot-20261011; GNU FDL 1.2 or later}
\begin{document}
\maketitle
\noindent\textit{Editorial note.} All delivered mathematical text is the author's own, copied verbatim from the pinned Stacks Project chapter file cohomology.tex; no statement or proof step was rewritten, reordered or removed.\par\medskip
\noindent\textit{Editorial note.} The counted span cites the lemmas recorded in assembly.external\_references with canonical tags as navigation pointers.\par\medskip
\noindent\textit{Allowed standard prerequisites (editorial).}\begin{itemize}
\item Evaluation map. Cohomology of Sheaves, Lemma cohomology-lemma-internal-hom-evaluate. The canonical map into the double dual.
\item Perfect case of the evaluation map. Cohomology of Sheaves, Lemma cohomology-lemma-internal-hom-evaluate-isom. The map is an isomorphism for perfect K.
\item Internal hom from strictly perfect. Cohomology of Sheaves, Lemma cohomology-lemma-Rhom-strictly-perfect. K\textasciicircum{} or  represented by termwise duals.
\item Evaluation map from the dual. Cohomology of Sheaves, Lemma cohomology-lemma-dual. The map (equation-eval).
\end{itemize}\medskip
\ifdefined\hypertarget\hypertarget{sourcepointer-1}{}\fi
\noindent\textit{Source reference.} Original lemma 0A8U: sources/stacks-chow/tags.txt, source lines 12927--12927. Cross-chapter standard prerequisite cited by the author inside the counted span; it is used as a pointer and its own proof is not part of this item.\par\medskip
\ifdefined\hypertarget\hypertarget{sourcepointer-2}{}\fi
\noindent\textit{Source reference.} Original lemma 0G40: sources/stacks-chow/tags.txt, source lines 20110--20110. Cross-chapter standard prerequisite cited by the author inside the counted span; it is used as a pointer and its own proof is not part of this item.\par\medskip
\ifdefined\hypertarget\hypertarget{sourcepointer-3}{}\fi
\noindent\textit{Source reference.} Original lemma 08DL: sources/stacks-chow/tags.txt, source lines 10641--10641. Cross-chapter standard prerequisite cited by the author inside the counted span; it is used as a pointer and its own proof is not part of this item.\par\medskip
\ifdefined\hypertarget\hypertarget{sourcepointer-4}{}\fi
\noindent\textit{Source reference.} Original lemma 08I1: sources/stacks-chow/tags.txt, source lines 10799--10799. Cross-chapter standard prerequisite cited by the author inside the counted span; it is used as a pointer and its own proof is not part of this item.\par\medskip
\ifdefined\hypertarget\hypertarget{sourcepointer-5}{}\fi
\noindent\textit{Source reference.} Original equation 08I2: sources/stacks-chow/tags.txt, source lines 10800--10800. Cross-chapter standard prerequisite cited by the author inside the counted span; it is used as a pointer and its own proof is not part of this item.\par\medskip
\ifdefined\hypertarget\hypertarget{sourcepointer-6}{}\fi
\noindent\textit{Source reference.} Original lemma 08DM: sources/stacks-chow/tags.txt, source lines 10642--10642. Cross-chapter standard prerequisite cited by the author inside the counted span; it is used as a pointer and its own proof is not part of this item.\par\medskip
\begin{lemma}
\label{lemma-dual-perfect-complex}
Let $(X, \mathcal{O}_X)$ be a ringed space. Let $K$ be a perfect object of
$D(\mathcal{O}_X)$. Then $K^\vee = R\SheafHom(K, \mathcal{O}_X)$ is a
perfect object too and $(K^\vee)^\vee \cong K$. There are
functorial isomorphisms
$$
M \otimes^\mathbf{L}_{\mathcal{O}_X} K^\vee = R\SheafHom(K, M)
$$
and
$$
H^0(X, M \otimes^\mathbf{L}_{\mathcal{O}_X} K^\vee) =
\Hom_{D(\mathcal{O}_X)}(K, M)
$$
for $M$ in $D(\mathcal{O}_X)$.
\end{lemma}

\begin{proof}
By Lemma \ref{lemma-internal-hom-evaluate} there is a canonical map
$$
K = R\SheafHom(\mathcal{O}_X, \mathcal{O}_X)
\otimes_{\mathcal{O}_X}^\mathbf{L} K \longrightarrow
R\SheafHom(R\SheafHom(K, \mathcal{O}_X), \mathcal{O}_X) =
(K^\vee)^\vee
$$
which is an isomorphism by Lemma \ref{lemma-internal-hom-evaluate-isom}.
To check the other statements we will use without further mention that
formation of internal hom commutes with restriction to opens
(Lemma \ref{lemma-restriction-RHom-to-U}).
We may check $K^\vee$ is perfect locally on $X$.
By Lemma \ref{lemma-dual}
to see the final statement it suffices to check that the map
(\ref{equation-eval})
$$
M \otimes^\mathbf{L}_{\mathcal{O}_X} K^\vee
\longrightarrow
R\SheafHom(K, M)
$$
is an isomorphism. This is local on $X$ as well.
Hence it suffices to prove these two statements $K$ is represented
by a strictly perfect complex.

\medskip\noindent
Assume $K$ is represented by the strictly perfect complex
$\mathcal{E}^\bullet$. Then it follows from
Lemma \ref{lemma-Rhom-strictly-perfect}
that $K^\vee$ is represented by the complex whose terms are
$(\mathcal{E}^{-n})^\vee =
\SheafHom_{\mathcal{O}_X}(\mathcal{E}^{-n}, \mathcal{O}_X)$
in degree $n$. Since $\mathcal{E}^{-n}$ is a direct summand of a finite
free $\mathcal{O}_X$-module, so is $(\mathcal{E}^{-n})^\vee$.
Hence $K^\vee$ is represented by a strictly perfect complex too
and we see that $K^\vee$ is perfect.
To see that (\ref{equation-eval}) is an isomorphism, represent
$M$ by a complex $\mathcal{F}^\bullet$.
By Lemma \ref{lemma-Rhom-strictly-perfect} the complex
$R\SheafHom(K, M)$ is represented by the complex with terms
$$
\bigoplus\nolimits_{n = p + q}
\SheafHom_{\mathcal{O}_X}(\mathcal{E}^{-q}, \mathcal{F}^p)
$$
On the other hand, the object $M \otimes^\mathbf{L}_{\mathcal{O}_X} K^\vee$
is represented by the complex with terms
$$
\bigoplus\nolimits_{n = p + q}
\mathcal{F}^p \otimes_{\mathcal{O}_X} (\mathcal{E}^{-q})^\vee
$$
Thus the assertion that (\ref{equation-eval}) is an isomorphism
reduces to the assertion that the canonical map
$$
\mathcal{F}
\otimes_{\mathcal{O}_X}
\SheafHom_{\mathcal{O}_X}(\mathcal{E}, \mathcal{O}_X)
\longrightarrow
\SheafHom_{\mathcal{O}_X}(\mathcal{E}, \mathcal{F})
$$
is an isomorphism when $\mathcal{E}$ is a direct summand of a finite
free $\mathcal{O}_X$-module and $\mathcal{F}$ is any $\mathcal{O}_X$-module.
This follows immediately from the corresponding statement when
$\mathcal{E}$ is finite free.
\end{proof}

\end{document}
